\documentclass[a4paper]{article}

% \usepackage[default]{fontsetup}

\usepackage{fancyhdr}
\usepackage{extramarks}
\usepackage{amsmath}
\usepackage{amsthm}
\usepackage{amsfonts}
\usepackage{tikz}
\usepackage[plain]{algorithm}
\usepackage{algpseudocode}
\usepackage{enumerate}
\usepackage{tikz}

\usetikzlibrary{automata,positioning}

%
% Basic Document Settings
%  

\topmargin=-0.2in
\evensidemargin=0in
\oddsidemargin=0in
\textwidth=6.5in
\textheight=9.5in
\headsep=0.25in

\linespread{1.1}

\pagestyle{fancy}
\lhead{\hmwkAuthorName}
\chead{\hmwkClass : \hmwkTitle}
\rhead{\firstxmark}
\lfoot{\lastxmark}
\cfoot{\thepage}

\renewcommand\headrulewidth{0.4pt}
\renewcommand\footrulewidth{0.4pt}

\setlength\parindent{0pt}

%
% Create Problem Sections
%

\newcommand{\enterProblemHeader}[1]{
    \nobreak\extramarks{}{Problem \arabic{#1} continued on next page\ldots}\nobreak{}
    \nobreak\extramarks{Problem \arabic{#1} (continued)}{Problem \arabic{#1} continued on next page\ldots}\nobreak{}
}

\newcommand{\exitProblemHeader}[1]{
    \nobreak\extramarks{Problem \arabic{#1} (continued)}{Problem \arabic{#1} continued on next page\ldots}\nobreak{}
    \stepcounter{#1}
    \nobreak\extramarks{Problem \arabic{#1}}{}\nobreak{}
}

\newcommand*\circled[1]{\tikz[baseline=(char.base)]{
		\node[shape=circle,draw,inner sep=2pt] (char) {#1};}}


\setcounter{secnumdepth}{0}
\newcounter{partCounter}
\newcounter{homeworkProblemCounter}
\setcounter{homeworkProblemCounter}{1}
\nobreak\extramarks{Problem \arabic{homeworkProblemCounter}}{}\nobreak{}

%
% Homework Problem Environment
%
% This environment takes an optional argument. When given, it will adjust the
% problem counter. This is useful for when the problems given for your
% assignment aren't sequential. See the last 3 problems of this template for an
% example.
%

\newenvironment{homeworkProblem}[2][-1]{%
    \ifnum#1>0\relax
        \setcounter{homeworkProblemCounter}{#1}
    \fi
    \section{Problem \arabic{homeworkProblemCounter} (#2 points)}
    \setcounter{partCounter}{1}
    \enterProblemHeader{homeworkProblemCounter}
}{%
    \exitProblemHeader{homeworkProblemCounter}
}

%
% Homework Details
%   - Title
%   - Class
%   - Due date
%   - Name
%   - Student ID

\newcommand{\hmwkTitle}{Homework\ \#05}
\newcommand{\hmwkClass}{Probability \& Statistics for EECS}
\newcommand{\hmwkDueDate}{2026-10-25}
\newcommand{\hmwkAuthorName}{Jianxuan Yao}
\newcommand{\hmwkAuthorID}{2025521024}


%
% Title Page
%

\title{
    \vspace{2in}
    \textmd{\textbf{\hmwkClass:\\  \hmwkTitle}}\\
    \normalsize\vspace{0.1in}\small{Due\ on\ \hmwkDueDate\ at 23:59}\\
	\vspace{4in}
}

\author{
	Name: \textbf{\hmwkAuthorName} \\
	Student ID: \hmwkAuthorID}
\date{}

\renewcommand{\part}[1]{\textbf{\large Part \Alph{partCounter}}\stepcounter{partCounter}\\}

%
% Various Helper Commands
%

% Useful for algorithms
\newcommand{\alg}[1]{\textsc{\bfseries \footnotesize #1}}
% For derivatives
\newcommand{\deriv}[1]{\frac{\mathrm{d}}{\mathrm{d}x} (#1)}
% For partial derivatives
\newcommand{\pderiv}[2]{\frac{\partial}{\partial #1} (#2)}
% Integral dx
\newcommand{\dx}{\mathrm{d}x}
% Alias for the Solution section header
\newcommand{\solution}{\textbf{\large Solution}}
% Probability commands: Expectation, Variance, Covariance, Bias
\newcommand{\E}{\mathrm{E}}
\newcommand{\Var}{\mathrm{Var}}
\newcommand{\Cov}{\mathrm{Cov}}
\newcommand{\Bias}{\mathrm{Bias}}

\begin{document}


% \maketitle
% \thispagestyle{empty}
% \pagebreak

\date{
Due on Oct. 25, 2026, 23:59 UTC+8}
\title{SI 140A.02  Probability \& Statistics for EECS, Fall 2026 \\
Homework 5}
\maketitle
Read all the instructions below carefully before you start working on the assignment, and before you make a submission.
\begin{itemize}
    \item You are required to submit your homework in PDF format generated by \textbf{LaTeX}, otherwise you will get no points of this homework.
    \item You are required to write down all the major steps towards making your conclusions; otherwise you may obtain limited points of the problem.
    \item Write your homework in English; otherwise you will get no points of this homework.
    \item Any form of plagiarism will lead to $0$ point of this homework. 
\end{itemize}
\newpage

\begin{homeworkProblem}{20}

A/B testing is a randomized experiment used to compare two
treatments, Treatment A and Treatment B. Users arrive one by one
and are assigned to one of the treatments. Each trial is classified
as either a success or a failure.

Consider the following ``stay-with-a-winner'' procedure:
\begin{enumerate}[(i)]
    \item Assign the first user to Treatment A or Treatment B
    with equal probabilities.
    \item If the trial for the $n$th user is a success, assign
    the $(n+1)$st user to the same treatment. Otherwise, switch
    to the other treatment.
\end{enumerate}

The success probabilities for Treatments A and B are $a$ and $b$,
respectively, where $a\ne b$. Given the assigned treatment, these
success probabilities do not depend on the previous outcomes.

Let $p_n$ be the probability of success on the $n$th trial, and
let $a_n$ be the probability that Treatment A is assigned on the
$n$th trial.

\begin{enumerate}[(a)]
    \item Show that
    \[
        p_n=(a-b)a_n+b,
        \qquad
        a_{n+1}=(a+b-1)a_n+1-b.
    \]

    \item Use part (a) to show that
    \[
        p_{n+1}=(a+b-1)p_n+a+b-2ab.
    \]

    \item Assuming that the limit exists, find
    \[
        \lim_{n\to\infty}p_n.
    \]
\end{enumerate}

\end{homeworkProblem}

\solution

\begin{enumerate}
    \item[(a)] 
    
    $$
    a_1 = b_1 = \frac{1}{2}
    $$

    $$
    a_{n+1}
    = b_n (1-b) + a_n a
    $$

    $$
    b_{n+1}
    = a_n (1-a) + b_n b
    $$

    $$
    a_n + b_n = 1
    $$

    $$
    a_{n+1}
    = (1-a_n)(1-b) + a_n a
    = (a+b-1)a_n + 1-b
    $$

    $$
    p_n = a_n a + b_n b
    = a_n a + (1-a_n) b
    = (a-b)a_n + b
    $$

    \item[(b)]
    
    $$
    \begin{aligned}
    p_{n+1}
    &= a_{n+1} a + b_{n+1} b \\
    &= ((a+b-1)a_n + 1-b) a + ((a+b-1)b_n + 1-a) b \\
    &= (a+b-1)(a_n a + b_n b) + a + b - 2ab \\
    &= (a+b-1)p_n + a + b - 2ab
    \end{aligned}
    $$

    \item[(c)]
    
    Let $L = \lim_{n\to\infty}p_n$.

    $$
    L = (a+b-1)L + a + b - 2ab
    $$

    $$
    L = \frac{a + b - 2ab}{2 - a - b}
    $$
\end{enumerate}

\newpage

\begin{homeworkProblem}{20}

A fair coin is tossed repeatedly and independently, producing a
sequence of heads and tails. Let $N$ be the number of tosses
required to observe the first occurrence of the pattern $HH$.

Find the PMF of $N$.

\end{homeworkProblem}

\solution

The sequence's $n$-th position and $(n-1)$-th position must be $HH$, the $(n-2)$-th position must be $T$, and the first $n-3$ positions must not contain $HH$.

Let $a_k$ be the number of sequences of length $k$ that do not contain $HH$.

$$
a_1 = 2, \quad a_2 = 3
$$

$$
a_k = a_{k-1} + a_{k-2}
$$

$$
a_k = \frac{1}{\sqrt{5}} \left[ \left(\frac{1+\sqrt{5}}{2}\right)^{k+2} - \left(\frac{1-\sqrt{5}}{2}\right)^{k+2} \right]
$$

$$
P(N=n) = \frac{a_{n-3}}{2^n} = \frac{1}{2^n \sqrt{5}} \left[ \left(\frac{1+\sqrt{5}}{2}\right)^{n-1} - \left(\frac{1-\sqrt{5}}{2}\right)^{n-1} \right]
$$

\newpage


\begin{homeworkProblem}{15}

A system contains seven devices: four in Group A and three in
Group B. Each device fails independently with probability $p$,
where $0<p<1$.

Let $X$ be the number of failed devices in Group A, and let $T$
be the total number of failed devices in the system.

\begin{enumerate}[(a)]
    \item Find the PMF of $X$ and the PMF of $T$.

    \item Given that exactly three devices have failed, find
    the conditional PMF of $X$. Identify the distribution.

    \item Explain why the conditional distribution in part (b)
    does not depend on $p$.
\end{enumerate}

\end{homeworkProblem}

\solution

\begin{enumerate}
    \item[(a)]
    
    $$
    X \sim \text{Bin}(4, p)
    $$

    $$
    T \sim \text{Bin}(7, p)
    $$

    $$
    P(X=x) = \binom{4}{x} p^x (1-p)^{4-x}, \quad x \in \{0,1,2,3,4\}
    $$

    $$
    P(T=t) = \binom{7}{t} p^t (1-p)^{7-t}, \quad t \in \{0,1,2,3,4,5,6,7\}
    $$

    \item[(b)]

    Let $Y$ be the number of failed devices in Group B. Then $X + Y = T$.

    $$
    Y \sim \text{Bin}(3, p)
    $$

    $$
    P(X = k  \mid  T = 3)
    = \frac{P(X = k, T = 3)}{P(T = 3)}
    = \frac{P(X = k) P(Y = 3 - k)}{P(T = 3)}
    $$

    $$
    P(X = k  \mid  T = 3) = \frac{\binom{4}{k} \binom{3}{3-k}}{\binom{7}{3}}
    $$

    $$
    P(X = 0  \mid  T = 3) = \frac{\binom{4}{0}\binom{3}{3}}{35} = \frac{1 \times 1}{35} = \frac{1}{35}
    $$

    $$
    P(X = 1  \mid  T = 3) = \frac{\binom{4}{1}\binom{3}{2}}{35} = \frac{4 \times 3}{35} = \frac{12}{35}
    $$

    $$
    P(X = 2  \mid  T = 3) = \frac{\binom{4}{2}\binom{3}{1}}{35} = \frac{6 \times 3}{35} = \frac{18}{35}
    $$

    $$
    P(X = 3  \mid  T = 3) = \frac{\binom{4}{3}\binom{3}{0}}{35} = \frac{4 \times 1}{35} = \frac{4}{35}
    $$

    $$
    X \sim \text{HGeom}(7, 4, 3)
    $$

    \item[(c)]
    
    For every device, the probability of failure is same and independent. Given there are exactly three 
    failed devices, the three failed devices can be any three of the seven devices, and the probability of each combination is equal.

    Therefore, the problem is equivalent to randomly selecting three devices from the seven devices, and counting how many of them are from Group A. This is a hypergeometric distribution, which does not depend on $p$.
\end{enumerate}

\newpage

\begin{homeworkProblem}{15}

Two fair six-sided dice are rolled independently. Let $U$ and
$V$ be their respective outcomes, and define
\[
    D= \mid U-V \mid .
\]

\begin{enumerate}[(a)]
    \item Find the PMF of $D$.

    \item Find the CDF $F_D(t)$ for every real number $t$.

    \item Define $W=\min(D,2)$. Find the PMF of $W$.
\end{enumerate}

\end{homeworkProblem}

\solution

\begin{enumerate}
    \item[(a)]
    
    $$
    P(D=0) = \frac{1}{6}
    $$

    $$
    P(D=1) = \frac{10}{36} = \frac{5}{18}
    $$

    $$
    P(D=2) = \frac{8}{36} = \frac{2}{9}
    $$

    $$
    P(D=3) = \frac{6}{36} = \frac{1}{6}
    $$

    $$
    P(D=4) = \frac{4}{36} = \frac{1}{9}
    $$

    $$
    P(D=5) = \frac{2}{36} = \frac{1}{18}
    $$

    PMF of $D$:
    
    $$
    P(D=d) =
    \begin{cases}
        1/6,&d=0,\\
        5/18,&d=1,\\
        2/9,&d=2,\\
        1/6,&d=3,\\
        1/9,&d=4,\\
        1/18,&d=5,\\
        0,&\text{otherwise.}
    \end{cases}
    $$

    \item[(b)]
    
    $$
    F_D(t)
    =
    \begin{cases}
        0, & t < 0 \\
        \frac{1}{6}, & 0 \le t < 1 \\
        \frac{4}{9}, & 1 \le t < 2 \\
        \frac{2}{3}, & 2 \le t < 3 \\
        \frac{5}{6}, & 3 \le t < 4 \\
        \frac{17}{18}, & 4 \le t < 5 \\
        1, & t \ge 5
    \end{cases}
    $$

    \item[(c)]
    
    $$
    P(W = 0) = P(D = 0) = \frac{6}{36} = \frac{1}{6}
    $$

    $$
    P(W = 1) = P(D = 1) = \frac{10}{36} = \frac{5}{18}
    $$

    $$
    P(W = 2) = \frac{5}{9}
    $$

    $$
    p_W(w)
    =
    \begin{cases}
        \frac{1}{6}, & w = 0 \\
        \frac{5}{18}, & w = 1 \\
        \frac{5}{9}, & w = 2
    \end{cases}
    $$
\end{enumerate}

\newpage

\begin{homeworkProblem}{15}

There are two coins, A and B. Coin A lands heads with probability
$1/4$, while coin B lands heads with probability $3/4$.

One coin is selected with equal probabilities and then tossed
twice. The same selected coin is used for both tosses. Given
the selected coin, the two tosses are independent.

Let $C\in\{A,B\}$ denote the selected coin. Let $X$ and $Y$ be
the indicator random variables for heads on the first and
second tosses, respectively.

\begin{enumerate}[(a)]
    \item Find the PMF of $X$ and the PMF of $Y$.

    \item Are $X$ and $Y$ independent? Are they conditionally
    independent given $C$? Justify your answers.

    \item Given that both tosses are heads, find the
    conditional probability that coin B was selected.
\end{enumerate}

\end{homeworkProblem}

\solution

\begin{enumerate}
    \item[(a)]
    
    $$
    P(C=A) = P(C=B) = \frac{1}{2}
    $$

    $$
    P(X=0) = P(X=0 \mid C=A)P(C=A) + P(X=0 \mid C=B)P(C=B) = \frac{1}{2}
    $$
    
    $$
    P(X=1) = P(X=1 \mid C=A)P(C=A) + P(X=1 \mid C=B)P(C=B) = \frac{1}{2}
    $$

    $$
    P(Y=0) = P(Y=0 \mid C=A)P(C=A) + P(Y=0 \mid C=B)P(C=B) = \frac{1}{2}
    $$

    $$
    P(Y=1) = P(Y=1 \mid C=A)P(C=A) + P(Y=1 \mid C=B)P(C=B) = \frac{1}{2}
    $$

    \item[(b)]
    
    $X$ and $Y$ are not independent, and they are conditionally independent given $C$.

    $$
    P(X=1, Y=1) = P(X=1, Y=1 \mid C=A)P(C=A) + P(X=1, Y=1 \mid C=B)P(C=B) = \frac{5}{16}
    $$

    $$
    P(X=1)P(Y=1) = \frac{1}{4} \neq P(X=1, Y=1)
    $$

    And the question stem has already stated that given the selected coin, the two tosses are independent. As a result, $X$ and $Y$ are conditionally independent given $C$.

    \item[(c)]
    
    $$
    \begin{aligned}
    P(C=B \mid X=1, Y=1)
    &= \frac{P(X=1, Y=1 \mid C=B)P(C=B)}{P(X=1, Y=1)} \\
    &= \frac{\frac{9}{16} \cdot \frac{1}{2}}{\frac{5}{16}} \\
    &= \frac{9}{10}
    \end{aligned}
    $$
\end{enumerate}

\newpage


\begin{homeworkProblem}{15}

A random variable $X$ has PMF
\[
    \Pr(X=1)=\frac12,\qquad
    \Pr(X=2)=\frac14,\qquad
    \Pr(X=3)=\Pr(X=4)=\frac18.
\]
Let $U$ have the discrete uniform distribution on
$\{1,2,3,4\}$.

Throughout this problem, logarithms have base 2, and entropy
is measured in bits.

\begin{enumerate}[(a)]
    \item Find $H(X)$ and $H(U)$.

    \item Define
    \[
        Y=
        \begin{cases}
            1,&X\le2,\\
            0,&X>2,
        \end{cases}
        \qquad
        Z=
        \begin{cases}
            1,&X=1,\\
            0,&X\ne1.
        \end{cases}
    \]
    Find the PMF of each variable and calculate $H(Y)$
    and $H(Z)$.

    \item Is the size of the support sufficient to determine
    the entropy of a discrete random variable?
    Explain using distributions from this problem.
\end{enumerate}

\end{homeworkProblem}

\solution

\begin{enumerate}

    \item[(a)]
    
    $$
    H(X) = -\sum_{x=1}^{4} P(X=x) \log_2 P(X=x)
    = \frac{7}{4}
    $$

    $$
    H(U) = \log_2 4 = 2
    $$

    \item[(b)]
    
    $$
    P(Y=1) = P(X=1) + P(X=2) = \frac{3}{4}
    $$

    $$
    P(Y=0) = P(X=3) + P(X=4) = \frac{1}{4}
    $$

    $$
    P(Y=y) = 
    \begin{cases}
        \frac{3}{4}, & y=1 \\
        \frac{1}{4}, & y=0
    \end{cases}
    $$

    $$
    H(Y) = 2 - \frac{3}{4} \log_2 3
    $$

    $$
    P(Z=z) = 
    \begin{cases}
        \frac{1}{2}, & z=1 \\
        \frac{1}{2}, & z=0
    \end{cases}
    $$

    $$
    H(Z) = 1
    $$

    \item[(c)]
    
    The size of the support is not sufficient to determine the entropy of a discrete random variable.
    
    For example, $X$ and $U$ both have support size 4, but $H(X) = \frac{7}{4} < H(U) = 2$. Similarly, $Y$ and $Z$ both have support size 2, but $H(Y) = 2 - \frac{3}{4} \log_2 3 > H(Z) = 1$.
\end{enumerate}

\end{document}